%% chapter3.tex – Background and Technologies
%% Group No. 2 | Dept. of CSE (Data Science), MACE Kothamangalam
%% Team: Muhammed Ajmal P (MAC23CD040), Muhammed Adil PP (MAC23CD039),
%%       Mohammed Shadeed P (MAC23CD037), Sinan Rahman (LMAC23CD066)
%% Guide: Prof. Eldo P Elias | May 2026

\chapter{Background and Technologies}
\label{chap:background}

\section{Internet of Things Architecture}

The Internet of Things (IoT) refers to a network of physical objects embedded with
sensors, processing units, and communication interfaces that enable data collection
and exchange~\cite{alfuqaha2015iot}.
A typical IoT architecture comprises three layers:

\begin{enumerate}
  \item \textbf{Perception layer} – sensors and actuators that interact with the
        physical world (MPU6050, NEO-6M GPS, SIM800L).
  \item \textbf{Network layer} – communication protocols that transport data
        (WiFi 802.11~b/g/n, GSM/GPRS).
  \item \textbf{Application layer} – data processing, storage, and user interfaces
        (FastAPI backend, React dashboard).
\end{enumerate}

WaveGuard~V4 maps cleanly onto this three-layer model, with the ESP32 acting as the
integration hub between the perception and network layers.

\section{ESP32 Microcontroller}

The ESP32 DevKit~v1 is a 32-bit dual-core SoC manufactured by Espressif Systems,
featuring a Xtensa LX6 processor running at up to 240~MHz~\cite{espressif2023esp32}.
Relevant specifications for WaveGuard~V4 are listed in Table~\ref{tab:esp32spec}.

\begin{table}[H]
\centering
\caption{ESP32 DevKit v1 Key Specifications}
\label{tab:esp32spec}
\begin{tabular}{ll}
\toprule
\textbf{Parameter}   & \textbf{Value}        \\
\midrule
Architecture         & Xtensa LX6 dual-core  \\
Clock Speed          & 240~MHz               \\
SRAM                 & 520~KB                \\
Flash (on-board)     & 4~MB                  \\
Operating Voltage    & 3.3~V (I/O)           \\
WiFi                 & 802.11 b/g/n (2.4~GHz)\\
Bluetooth            & BT 4.2 / BLE          \\
GPIO Pins            & 34 (programmable)     \\
ADC                  & 12-bit, up to 18 ch.  \\
I2C Controllers      & 2                     \\
UART Controllers     & 3                     \\
Deep Sleep Current   & $\sim$10~µA            \\
\bottomrule
\end{tabular}
\end{table}

The ESP32's integrated WiFi stack allows it to serve an HTTP/WebSocket server directly,
eliminating the need for a separate WiFi module.
Its deep-sleep capability (10~µA) enables future duty-cycling to extend battery life.

\section{MPU6050 Inertial Measurement Unit}

The MPU6050~\cite{invensense2013mpu6050} is a 6-DOF IMU combining a three-axis
gyroscope and a three-axis accelerometer in a single 4~×~4~mm QFN package.
The device communicates via I2C at up to 400~kHz.

\subsection{Accelerometer Principles}

A MEMS accelerometer operates on the principle of a proof mass suspended by
spring-like poly-silicon structures.
An applied acceleration displaces the mass, changing the differential capacitance
between the fixed and moving comb structures.
This capacitance change is converted to a digital value via an on-chip ADC.

The sensitivity is configurable:
\begin{equation}
  a [\textrm{g}] = \frac{\textrm{Raw ADC count}}{S_f}
  \label{eq:accel_scale}
\end{equation}
where the scale factor $S_f$ is 16\,384~LSB/g for the $\pm$2~g range selected in
WaveGuard~V4.

\subsection{RMS-Based Wave Metric}

Raw acceleration $\mathbf{a}(t) = [a_x, a_y, a_z]$ is processed to extract the
vertical RMS acceleration over a 1-second window of $N = 50$ samples:

\begin{equation}
  a_{\rm rms} = \sqrt{\frac{1}{N}\sum_{i=1}^{N}\left[(a_{xi}-\bar{a}_x)^2 +
    (a_{yi}-\bar{a}_y)^2 + (a_{zi}-\bar{a}_z)^2\right]}
  \label{eq:rms}
\end{equation}

Gravitational offset is removed by subtracting the mean $\bar{a}$ over the calibration
window, leaving only dynamic (wave-induced) acceleration.
The MPU6050 register map for the $\pm$2~g range and the I2C address 0x68 are
used~\cite{invensense2013mpu6050}.

\section{SIM800L GSM/GPRS Module}

The SIM800L is a quad-band (850/900/1800/1900~MHz) GSM/GPRS module with an onboard
TCP/IP stack~\cite{simcom2015sim800l}.
WaveGuard~V4 uses only the SMS (AT+CMGS) facility.
Key electrical characteristics relevant to our power design:

\begin{itemize}
  \item Supply voltage: 3.7--4.2~V
  \item Idle current: $\sim$18~mA
  \item Transmit peak: up to 2~A (pulse, $\approx$500~ms)
  \item Typical SMS delivery time: 8--12~s from AT command issuance
\end{itemize}

The 2~A transmit spike necessitates a dedicated low-ESR capacitor bank (1\,000~µF
electrolytic + 100~µF ceramic) on the SIM800L power rail to prevent brownout of the
ESP32 during transmission.

\section{NEO-6M GPS Module}

The u-blox NEO-6M~\cite{ublox2011neo6m} is a compact GPS receiver producing NMEA~0183
sentences at 9\,600~baud over UART.
Position accuracy is approximately 2.5~m CEP under open-sky conditions.
The module acquires a first fix in approximately 30~s (cold start) and provides
1~Hz position updates thereafter.
WaveGuard~V4 reads \verb|$GPGGA| sentences to extract latitude/longitude and
attaches them to each sensor record stored in the database.

\section{WiFi Communication (IEEE 802.11 b/g/n)}

The ESP32's 802.11~b/g/n radio operates in the 2.4~GHz ISM band.
In \textit{access-point mode} the ESP32 creates its own network (SSID:
\texttt{WaveGuard\_AP}) to which a shore-side laptop connects for dashboard access.
In \textit{station mode} it associates with an existing router.
Typical line-of-sight range in open water is $\approx$50~m at 2.4~GHz; the
dashboard is therefore accessible from the shore station.

\section{Power Management}

\subsection{Li-ion Cell (18650)}

The 18650 format Li-ion cell provides a nominal 3.7~V with 2\,200~mAh capacity.
Usable energy is approximately:

\begin{equation}
  E = V_{\rm nom} \times C = 3.7~\textrm{V} \times 2.2~\textrm{Ah}
    = 8.14~\textrm{Wh}
\end{equation}

At an average system draw of $\approx$150~mA (active mode), the theoretical runtime is:

\begin{equation}
  t = \frac{2200~\textrm{mAh}}{150~\textrm{mA}} \approx 14.7~\textrm{h}
\end{equation}

Accounting for SMS-burst peaks and the boost converter efficiency ($\eta \approx
88\%$), empirically measured runtime is 6--8~hours.

\subsection{MT3608 Boost Converter}

The MT3608 is a 2~MHz fixed-frequency PWM step-up converter capable of supplying
2~A from a cell voltage of 2.0--24~V~\cite{aerosemi2010mt3608}.
It is used to generate the 5~V rail that powers the ESP32 (via its onboard LDO)
and ancillary 5~V logic.

\subsection{TP4056 Charger}

The TP4056 linear charger IC manages safe CC/CV charging of the 18650 cell at 1~A
from a USB~5~V source~\cite{toppower2012tp4056}.
A protection circuit module (DW01A + FS8205A) is included to guard against
over-discharge and over-voltage.

\section{I2C Serial Protocol}

I2C is a two-wire master–slave protocol~\cite{nxp2021i2c}.
The ESP32 acts as master on GPIO21 (SDA) and GPIO22 (SCL) with 4.7~k$\Omega$
pull-up resistors to 3.3~V.
The MPU6050 is assigned I2C address \texttt{0x68} (AD0 pin tied to GND).
The bus runs at 400~kHz (fast mode), allowing 50~Hz sensor reads with margin.

\section{Real-Time Alert Logic}

Algorithm~\ref{alg:alert} describes the onboard alert decision loop:

\begin{lstlisting}[language=C++, caption={WaveGuard V4 Alert Decision Loop (pseudocode)}, label={alg:alert}]
void loop() {
  float rms = computeRMS();        // Equation 3.2 over 50 samples

  String level;
  if      (rms < 0.03) level = "CALM";
  else if (rms < 0.075) level = "MODERATE";
  else                  level = "HIGH_WAVE";

  postToBackend(rms, level, gpsLat, gpsLon);

  if (level == "HIGH_WAVE" && smsEnabled) {
    sendSMS(alertNumber, buildSMSMessage(rms, gpsLat, gpsLon));
  }

  delay(1000);   // 1-second sampling interval
}
\end{lstlisting}
